\documentclass[11pt]{article}
\usepackage[margin=1in]{geometry}
\usepackage[T1]{fontenc}
\usepackage[utf8]{inputenc}
\usepackage{amsmath,amssymb,mathtools}
\usepackage{booktabs}
\usepackage{hyperref}
\usepackage{xcolor}
\usepackage{listings}

\hypersetup{
  colorlinks=true,
  linkcolor=blue!60!black,
  urlcolor=blue!60!black
}

\definecolor{codebg}{RGB}{248,248,248}
\definecolor{codeborder}{RGB}{210,210,210}
\definecolor{codecomment}{RGB}{80,120,80}
\definecolor{codekeyword}{RGB}{30,60,150}

\lstdefinestyle{reportcode}{
  backgroundcolor=\color{codebg},
  basicstyle=\ttfamily\small,
  breaklines=true,
  columns=fullflexible,
  commentstyle=\color{codecomment},
  frame=single,
  framerule=0.5pt,
  keywordstyle=\color{codekeyword},
  rulecolor=\color{codeborder},
  keepspaces=true,
  showstringspaces=false,
  tabsize=2
}

\lstset{style=reportcode, language=Python}

\setlength{\parindent}{0pt}
\setlength{\parskip}{0.5em}

\newcommand{\R}{\mathbb{R}}
\newcommand{\ysp}{y_{\mathrm{sp}}}
\newcommand{\uvec}{u}
\newcommand{\xvec}{x}
\newcommand{\Aaug}{A_{\mathrm{aug}}}
\newcommand{\Baug}{B_{\mathrm{aug}}}
\newcommand{\Caug}{C_{\mathrm{aug}}}
\newcommand{\Qout}{Q_{\mathrm{out}}}
\newcommand{\Rin}{R_{\mathrm{in}}}
\newcommand{\Hp}{H_p}
\newcommand{\Hc}{H_c}
\newcommand{\argmin}{\mathop{\mathrm{arg\,min}}}

\usepackage{graphicx}
\usepackage{float}
\graphicspath{{figures/weekly_meeting_2026_04_17_v2/}}

\title{Weekly Meeting Summary: Disturbance-Mode RL-MPC Results}
\author{Prepared for supervisor meeting}
\date{2026-04-17}

\begin{document}
\maketitle

\section*{Scope}

This report is restricted to disturbance-mode results only. Combined-agent,
matrix, structured-matrix, and residual notebooks are excluded. The only
model-update family kept here is reidentification.

The notebook families covered are:
\begin{itemize}
\item distillation:
      \texttt{distillation\_RL\_assisted\_MPC\_horizons\_unified.ipynb},
      \texttt{distillation\_RL\_assisted\_MPC\_horizons\_dueling\_unified.ipynb},
      \texttt{distillation\_RL\_assisted\_MPC\_weights\_unified.ipynb},
      \texttt{distillation\_RL\_assisted\_MPC\_reidentification\_unified.ipynb}
\item polymer:
      \texttt{RL\_assisted\_MPC\_horizons\_unified.ipynb},
      \texttt{RL\_assisted\_MPC\_horizons\_dueling\_unified.ipynb},
      \texttt{RL\_assisted\_MPC\_weights\_unified.ipynb},
      \texttt{RL\_assisted\_MPC\_reidentification\_unified.ipynb}
\end{itemize}

For each included method, the figures focus on:
\begin{itemize}
\item average reward comparison against the paired MPC bundle
\item last-episode MPC comparison
\end{itemize}

\section*{Selection Notes}

\begin{itemize}
\item Distillation is presented first, per your request.
\item The latest stored April 16, 2026 distillation horizon, dueling-horizon,
      and weights runs use the newer unified reward profile introduced on
      April 16, 2026.
\item The latest stored disturbance reidentification result for distillation is
      still the TD3 run from April 16, 2026.
\item On polymer, I found qualifying disturbance bundles for DQN horizon,
      dueling horizon, weights, and reidentification.
\end{itemize}

\section*{Executive Summary}

\begin{itemize}
\item Distillation horizon adaptation is currently the strongest family in the
      report. The April 16, 2026 DQN and dueling runs both beat the paired MPC
      baseline, and the dueling variant is more concentrated around a small set
      of horizon recipes.
\item Distillation SAC weights are also positive on the newer reward profile,
      with final gap \(+1.520\) and last-50 gap \(+2.272\).
\item Distillation reidentification is still the most mixed result: strong
      mid-run collapse followed by recovery to a positive final gap \(+0.060\).
\item Polymer dueling horizon is still the best polymer disturbance result on
      disk. It beats the paired MPC baseline in 190 of 200 episodes and ends
      with final gap \(+1.692\).
\item Polymer weights are also clearly positive, with final gap \(+1.515\) and
      last-50 gap \(+1.732\).
\item Polymer DQN horizon remains positive and clearly beats MPC, but it is now
      weaker than the latest polymer dueling and weights runs.
\item Polymer reidentification is still slightly negative at the end, with final
      gap \(-0.070\), so polymer remains a weak benchmark for evaluating the
      model-update idea.
\end{itemize}

\section*{Distillation}

\begin{table}[H]
\centering
\small
\begin{tabular}{lrrrcr}
\toprule
Method & Agent & RL Final & MPC Final & Final Gap & Better Eps \\
\midrule
DQN horizon & DDQN & 19.433 & 16.079 & +3.354 & 118/200 \\
Dueling horizon & Dueling DDQN & 18.914 & 16.079 & +2.835 & 186/200 \\
Weights & SAC & 17.599 & 16.079 & +1.520 & 169/200 \\
Reidentification & TD3 & -0.205 & -0.265 & +0.060 & 70/200 \\
\bottomrule
\end{tabular}
\caption{Latest stored distillation disturbance runs with paired MPC comparison.}
\end{table}

\subsection*{DQN Horizon}

The latest stored disturbance DQN horizon run is
\texttt{distillation\_horizon\_disturb\_fluctuation\_standard\_unified/20260416\_192434},
compared against
\texttt{distillation\_compare\_horizon\_disturb\_fluctuation\_standard/20260416\_192443}.

\begin{itemize}
\item Final reward is 19.433 versus paired MPC 16.079, so the final gap is
      \(+3.354\). The last-50 gap is still positive at \(+0.299\).
\item The nominal MPC pair \((6, 3)\) is the most used recipe, but only at
      7.88\% of logged selections. The policy remains broad rather than
      collapsing to one horizon pair.
\item The best reward occurs at episode 131, but the run still has large
      negative excursions, with worst gap \(-11.253\) at episode 51.
\end{itemize}

\begin{figure}[H]
\centering
\includegraphics[width=0.49\textwidth]{../Distillation/Results/distillation_compare_horizon_disturb_fluctuation_standard/20260416_192443/compare_rewards.png}
\hfill
\includegraphics[width=0.49\textwidth]{../Distillation/Results/distillation_compare_horizon_disturb_fluctuation_standard/20260416_192443/compare_outputs_last_episode.png}
\caption{Distillation DQN horizon: average reward comparison, last-episode MPC
comparison.}
\end{figure}

\subsection*{Dueling Horizon}

The latest stored disturbance dueling-horizon run is
\texttt{distillation\_dueling\_horizon\_disturb\_fluctuation\_standard\_unified/20260416\_192302},
compared against
\texttt{distillation\_compare\_dueling\_horizon\_disturb\_fluctuation\_standard/20260416\_192310}.

\begin{itemize}
\item Final reward is 18.914 versus paired MPC 16.079, for final gap
      \(+2.835\). The last-50 gap is much stronger at \(+2.720\), and RL is
      better in 186 of 200 episodes.
\item The policy is much more concentrated than plain DQN horizon:
      \((6, 3)\) accounts for 53.36\% of logged selections, followed by
      \((8, 2)\) at 20.27\% and \((13, 5)\) at 16.12\%.
\item This currently looks like the more stable distillation horizon method:
      smaller final gap than DQN, but a much stronger late-run gap.
\end{itemize}

\begin{figure}[H]
\centering
\includegraphics[width=0.49\textwidth]{../Distillation/Results/distillation_compare_dueling_horizon_disturb_fluctuation_standard/20260416_192310/compare_rewards.png}
\hfill
\includegraphics[width=0.49\textwidth]{../Distillation/Results/distillation_compare_dueling_horizon_disturb_fluctuation_standard/20260416_192310/compare_outputs_last_episode.png}
\caption{Distillation dueling horizon: average reward comparison, last-episode
MPC comparison.}
\end{figure}

\subsection*{Weights}

The latest stored disturbance weights run is the SAC result
\texttt{distillation\_weights\_sac\_disturb\_fluctuation\_standard\_unified/20260416\_192555},
compared against
\texttt{distillation\_compare\_weights\_sac\_disturb\_fluctuation\_standard/20260416\_192602}.

\begin{itemize}
\item Final reward is 17.599 versus paired MPC 16.079, so the final gap is
      \(+1.520\). The last-50 gap is larger at \(+2.272\), which means the
      late part of training is stronger than the final episode alone suggests.
\item The four logged penalty multipliers average
      \([1.146,\ 1.162,\ 0.999,\ 1.033]\) and end at
      \([0.837,\ 1.100,\ 0.761,\ 1.014]\).
\item This is a useful positive result on the new reward profile even though the
      latest qualifying run is SAC rather than TD3.
\end{itemize}

\begin{figure}[H]
\centering
\includegraphics[width=0.49\textwidth]{../Distillation/Results/distillation_compare_weights_sac_disturb_fluctuation_standard/20260416_192602/compare_rewards.png}
\hfill
\includegraphics[width=0.49\textwidth]{../Distillation/Results/distillation_compare_weights_sac_disturb_fluctuation_standard/20260416_192602/compare_outputs_last_episode.png}
\caption{Distillation weights: average reward comparison, last-episode MPC
comparison.}
\end{figure}

\subsection*{Reidentification}

The latest stored disturbance reidentification run is
\texttt{distillation\_reidentification\_td3\_disturb\_fluctuation\_mismatch\_unified/20260416\_113030},
compared against
\texttt{distillation\_compare\_reidentification\_td3\_disturb\_fluctuation\_mismatch/20260416\_113035}.

\begin{itemize}
\item Final reward is \(-0.205\) versus paired MPC \(-0.265\), for final gap
      \(+0.060\). The last-50 gap is \(+0.025\).
\item The main issue is still the mid-run collapse. The worst gap is
      \(-4.237\) at episode 51, after which the run recovers and finally reaches
      its best reward at episode 200.
\item The blend factors stay moderate overall:
      \(\eta_A \approx 0.323\) mean and 0.385 last,
      \(\eta_B \approx 0.355\) mean and 0.398 last. Mean coefficient clipping is
      only 2.58\%, much lower than the polymer reidentification case.
\end{itemize}

\begin{figure}[H]
\centering
\includegraphics[width=0.49\textwidth]{../Distillation/Results/distillation_compare_reidentification_td3_disturb_fluctuation_mismatch/20260416_113035/compare_rewards.png}
\hfill
\includegraphics[width=0.49\textwidth]{../Distillation/Results/distillation_compare_reidentification_td3_disturb_fluctuation_mismatch/20260416_113035/compare_outputs_last_episode.png}
\caption{Distillation reidentification: average reward comparison, last-episode
MPC comparison.}
\end{figure}

\section*{Polymer}

\begin{table}[H]
\centering
\small
\begin{tabular}{lrrrcr}
\toprule
Method & Agent & RL Final & MPC Final & Final Gap & Better Eps \\
\midrule
DQN horizon & DDQN & -2.533 & -3.843 & +1.310 & 196/200 \\
Dueling horizon & Dueling DDQN & -2.726 & -4.417 & +1.692 & 190/200 \\
Weights & TD3 & -2.903 & -4.417 & +1.515 & 195/200 \\
Reidentification & TD3 & -4.487 & -4.417 & -0.070 & 10/200 \\
\bottomrule
\end{tabular}
\caption{Latest stored polymer disturbance runs with paired MPC comparison.}
\end{table}

\subsection*{DQN Horizon}

This notebook changes the MPC prediction and control horizons online with a DQN
policy. The latest stored disturbance run is
\texttt{horizon\_disturb\_unified/20260407\_175020}, compared against
\texttt{disturb\_compare\_horizon\_unified/20260407\_175031}.

\begin{itemize}
\item Final reward is \(-2.533\) versus paired MPC \(-3.843\), so the final gap
      is \(+1.310\). The last-50 gap is \(+1.213\), and RL is better in
      196 of 200 episodes.
\item The nominal MPC pair \((9, 3)\) is still the most used recipe, but only
      at 3.22\% of logged horizon selections. The policy is spread broadly over
      many longer-horizon combinations.
\item The best reward occurs late at episode 197. This remains a strong polymer
      disturbance result, but it is now behind the latest dueling and weights
      runs.
\end{itemize}

\begin{figure}[H]
\centering
\includegraphics[width=0.49\textwidth]{../Polymer/Results/disturb_compare_horizon_unified/20260407_175031/compare_rewards.png}
\hfill
\includegraphics[width=0.49\textwidth]{../Polymer/Results/disturb_compare_horizon_unified/20260407_175031/compare_outputs_last_episode.png}
\caption{Polymer DQN horizon: average reward comparison, last-episode MPC
comparison.}
\end{figure}

\subsection*{Dueling Horizon}

This notebook uses the dueling DDQN horizon supervisor. The latest stored
disturbance run is
\texttt{dueling\_horizon\_disturb\_unified/20260417\_035217}, compared against
\texttt{disturb\_compare\_dueling\_horizon\_unified/20260417\_035231}.

\begin{itemize}
\item Final reward is \(-2.726\) versus paired MPC \(-4.417\), so the final gap
      is \(+1.692\). The last-50 gap is \(+1.745\), and RL is better in
      190 of 200 episodes.
\item Unlike the distillation dueling case, polymer does not concentrate on the
      nominal MPC pair \((9, 3)\). The nominal pair is only 5.79\% of logged
      horizon selections and 0.83\% of decision points.
\item The most-used horizon recipes are all non-nominal, led by \((19, 5)\),
      \((9, 8)\), and \((17, 4)\). This is still the strongest polymer
      disturbance result currently stored.
\end{itemize}

\begin{figure}[H]
\centering
\includegraphics[width=0.49\textwidth]{../Polymer/Results/disturb_compare_dueling_horizon_unified/20260417_035231/compare_rewards.png}
\hfill
\includegraphics[width=0.49\textwidth]{../Polymer/Results/disturb_compare_dueling_horizon_unified/20260417_035231/compare_outputs_last_episode.png}
\caption{Polymer dueling horizon: average reward comparison, last-episode MPC
comparison.}
\end{figure}

\subsection*{Weights}

This notebook lets TD3 scale the output and input penalty weights online. The
latest stored disturbance run is
\texttt{td3\_weights\_disturb/20260417\_031734}, compared against
\texttt{disturb\_compare\_td3\_weights/20260417\_031746}.

\begin{itemize}
\item Final reward is \(-2.903\) versus paired MPC \(-4.417\), so the final gap
      is \(+1.515\). The last-50 gap is \(+1.732\), and RL is better in
      195 of 200 episodes.
\item The four logged penalty multipliers average
      \([1.403,\ 1.495,\ 1.263,\ 1.333]\) and end at
      \([1.957,\ 1.703,\ 1.176,\ 1.246]\).
\item This is a clear positive polymer result, but it still trails the latest
      dueling-horizon run on final reward gap.
\end{itemize}

\begin{figure}[H]
\centering
\includegraphics[width=0.49\textwidth]{../Polymer/Results/disturb_compare_td3_weights/20260417_031746/compare_rewards.png}
\hfill
\includegraphics[width=0.49\textwidth]{../Polymer/Results/disturb_compare_td3_weights/20260417_031746/compare_outputs_last_episode.png}
\caption{Polymer weights: average reward comparison, last-episode MPC
comparison.}
\end{figure}

\subsection*{Reidentification}

This notebook identifies low-rank model corrections online and lets TD3 blend
the nominal and identified models. The latest stored disturbance run is
\texttt{td3\_reidentification\_disturb/20260415\_120803}, compared against
\texttt{disturb\_compare\_td3\_reidentification/20260415\_120810}.

\begin{itemize}
\item Final reward is \(-4.487\) versus paired MPC \(-4.417\), so the final gap
      is \(-0.070\). The last-50 gap is \(-0.087\), and RL is better in only
      10 of 200 episodes.
\item The mean blend factors are \(\eta_A \approx 0.421\) and
      \(\eta_B \approx 0.411\), but by the end the policy shifts strongly toward
      \(B\): \(\eta_A = 0.108\), \(\eta_B = 0.846\).
\item Mean coefficient clipping fraction is 12.18\%, much higher than the
      distillation reidentification run. Polymer is still not giving a useful
      positive signal for this family.
\end{itemize}

\begin{figure}[H]
\centering
\includegraphics[width=0.49\textwidth]{../Polymer/Results/disturb_compare_td3_reidentification/20260415_120810/compare_rewards.png}
\hfill
\includegraphics[width=0.49\textwidth]{../Polymer/Results/disturb_compare_td3_reidentification/20260415_120810/compare_outputs_last_episode.png}
\caption{Polymer reidentification: average reward comparison, last-episode MPC
comparison.}
\end{figure}

\section*{Van de Vusse Motivation And Current Stage}

The current motivation for Van de Vusse is still the same: polymer is giving a
good signal for horizon and weight adaptation, but it is still a weak benchmark
for the model-update idea. Polymer reidentification remains slightly worse than
MPC, and the removed matrix-style families were also not giving strong gains.
That makes it hard to judge whether the adaptation idea itself is weak or
whether the polymer benchmark is simply not sensitive enough.

Van de Vusse is therefore being developed as a benchmark where operating-point
changes should create a clearer need for model adaptation. The current stage is
still pre-RL:
\begin{itemize}
\item nonlinear plant added under \texttt{systems/vandevusse/}
\item direct-local-linearization system identification in place
\item corrected benchmark feed-temperature handling
\item offset-free MPC baseline notebook running in disturbance mode with the
      canonical \texttt{ca0\_blocks} disturbance
\item observer-update cleanup and observer-pole search utilities added
\item no Van de Vusse RL result bundle yet, so it is not part of the tables
\end{itemize}

\end{document}
